\documentclass[fontset=none]{ctexart}
\usepackage{qp-m3}
\begin{document}
\renewcommand{\qpzhangming}{第九章\quad 静电场及其应用}
\renewcommand{\qpjianming}{检测·试产片2}
\keshi{9.3 电场 电场强度（检测场·A5 分档举证）}
\begin{multicols}{2}
\emergencystretch=1em

\tihao{1}\zu{短式×一行四连排} \tieside{简单}电场中某点的电场强度\nobreak(\kongwei)

\lxoptq{A．\(E/4\)；}{B．\(E/2\)；}{C．\(2E\)；}{D．\(4E\)}

\begin{ansblock}[A试-2-jc-01]
\ansitem{1}{D．}
\end{ansblock}

\tihao{2}\zu{短词×2×2 档} \tieside{简单}放入试探电荷后该点场强\nobreak(\kongwei)

\lxoptq{A．逐渐变大；}{B．逐渐变小；}{C．保持不变；}{D．无法判断}

\begin{ansblock}[A试-2-jc-02]
\ansitem{2}{C．}
\end{ansblock}

\tihao{3}\zu{长句×独立成行档} \tieside{中档}关于电场强度，下列说法正确的是\nobreak(\kongwei)

\lxoptq{A．电场中某点的电场强度与放在该点的试探电荷的电荷量成反比；}{B．电场强度方向与负试探电荷所受静电力方向相同；}{C．电场强度是矢量，其方向与正电荷在该点所受静电力的方向相同；}{D．若不放试探电荷，则该点电场强度为零}

\begin{ansblock}[A试-2-jc-03]
\ansitem{3}{C．}
\ansnote{详解}{场强由电场本身决定，与试探电荷有无、电荷量均无关（A、D 错）；方向规定沿正电荷受力方向（B 错）。}
\end{ansblock}

\end{multicols}

\end{document}
